% -*- coding: utf-8 -*-
% !TEX program = xelatex
\documentclass[plain,noanswer,medmath]{../../template/jnuexam}
\usepackage{../../template/kysx}

\begin{document}


\examtitle{2001}{4}


\loadproblems{3}{2001P1}%载入试卷一
\loadproblems{3}{2001P3}%载入试卷三


\makepart{填空题}{本题共5小题，每小题3分，满分15分}


\useproblem{3}{1}{1}%【同试卷三第一(1)题】


\begin{problem}
设$z = \e^{- x} - f(x - 2y)$，且当$y = 0$时，$z = x^2$，则$\frac{\pdz}{\pdx} =$ \fillin{}．
\end{problem}


\begin{problem}
设行列式$D = \begin{vmatrix}
  3 &  0 &  4 & 0 \\
  2 &  2 &  2 & 2 \\
  0 & -7 &  0 & 0 \\
  5 &  3 & -2 & 2
\end{vmatrix}$，则第$4$行各元素余子式之和的值为 \fillin{}．
\end{problem}


\useproblem{3}{1}{3}%【同试卷三第一(3)题】


\useproblem{3}{1}{4}%【同试卷三第一(4)题】


\makepart{选择题}{本题共5小题，每小题3分，满分15分}


\useproblem{3}{2}{1}%【同试卷三第二(1)题】


\useproblem{3}{2}{2}%【同试卷三第二(2)题】


\useproblem{3}{2}{3}%【同试卷三第二(3)题】


\begin{problem}
对于任意二事件$A$和$B$，与$A \cup B = B$不等价的是\pickout{}
\begin{abcd}
\item $A \subset B$．
\item $\bar B \subset \bar A$．
\item $A\bar B = \emptyset$．
\item $\bar AB = \emptyset$．
\end{abcd}
\end{problem}


\useproblem{1}{2}{5}%【同试卷一第二(5)题】


\makepart{}{本题满分5分}%【同试卷三第三题】

\useproblem{3}{3}{0}


\makepart{}{本题满分6分}%【同试卷三第四题】

\useproblem{3}{4}{0}


\makepart{}{本题满分6分}%【同试卷三第五题】

\useproblem{3}{5}{0}


\makepart{}{本题满分7分}

\begin{problem}
某商品进价为$a$（元/件），根据以往经验，当销售价为$b$（元/件）时，销售量为$c$件（$a$, $b$, $c$均为正常数，
且$b \ge \frac{4}{3}a$）．市场调查表明，销售价每下降$10\%$，销售量可增加$40\%$．现决定一次性降价，
试问，当销售价定为多少时，可获得最大利润？并求出最大利润．
\end{problem}


\makepart{}{本题满分6分}

\begin{problem}
设$f(x)$在区间$[0.1]$上连续，在$(0,1)$内可导，且满足$f(1) = 3\int_0^{1/3} \e^{1 - x^2} f(x)\dx $．
证明至少存在一点$\xi \in(0,1)$，使得$f'(\xi) = 2\xi f(\xi)$．
\end{problem}


\makepart{}{本题满分6分}

\begin{problem}
设函数$f(x)$在$(0, +\infty)$内连续，$f(1) = \frac{5}{2}$，且对所有$x,t \in(0, +\infty)$，满足条件
$$\int_1^{xt} f(u)\du = t\int_1^x f(u)\du + x\int_1^t f(u)\du .$$
求$f(x)$．
\end{problem}


\makepart{}{本题满分9分}%【同试卷三第九题】

\useproblem{3}{9}{0}


\makepart{}{本题满分8分}

\begin{problem}
设$\alpha_i = (a_{i1},a_{i2}, \cdots ,a_{in})^T$（$i = 1,2, \cdots ,r$；$r < n$）是$n$维实向量，
且$\alpha_1$,$\alpha_2$,$\cdots$,$\alpha_r$线性无关．已知$\beta = (b_1,b_2, \cdots ,b_n)^T$是线性方程组
$$\begin{cases}
  a_{11}x_1 + a_{12}x_2 + \cdots + a_{1n}x_n = 0, \\
  a_{21}x_1 + a_{22}x_2 + \cdots + a_{2n}x_n = 0, \\
  \cdots \cdots \\
  a_{r1}x_1 + a_{r2}x_2 + \cdots + a_{rn}x_n = 0,
\end{cases}$$
的非零解向量．试判断向量组$\alpha_1$,$\alpha_2$,$\cdots$,$\alpha_r$,$\beta$的线性相关性．
\end{problem}


\makepart{}{本题满分8分}%【同试卷三第十一题】

\useproblem{3}{11}{0}


\makepart{}{本题满分8分}

\begin{problem}
设随机变量$X$和$Y$的联合分布在以点$(0,1)$,$(1,0)$,$(1,1)$为顶点的三角形区域上服从均匀分布，
试求随机变量$U = X + Y$的方差．
\end{problem}


\end{document}
